\begin{table}[H]\centering
\caption{\textbf{Missingness of the cesarean-timing indicator is balanced across the week}}
\label{tab:timing_missing}
\small
\small\setlength{\tabcolsep}{4pt}\resizebox{\ifdim\width>\linewidth \linewidth\else\width\fi}{!}{%
\begin{tabular}{lcc}
\toprule
Sector & Weekday (\%) & Weekend (\%) \\
\midrule
For-profit & 6.81 & 6.71 \\
Public & 10.34 & 10.31 \\
\bottomrule
\end{tabular}
}
\begin{minipage}{\linewidth}\footnotesize
\textit{Notes:} Share of cesareans with a missing before/during-labor code, SINASC 2012--2024, the years the timing analyses use. The near-identical weekday and weekend rates imply the sector-level prelabor/in-labor split is not driven by differential missingness across the week. The balance does not carry over to Robson group 1, which receives the cesareans whose code is blank or ignored: there 11.2 percent of for-profit cesareans lack the code on weekdays and 9.8 percent on weekends (2014--2024), so part of that group's weekend dip may be uncoded prelabor cesareans (Table~\ref{tab:robson_validation}).
\end{minipage}
\end{table}
